\begin{textsample}{POS Dim 1 – sc – Score 49.00 – sc/sc\_inf\_e\_ef\_26.txt}  \label{ex:f1_pos_010}
3 ALFABETIZAÇÃO 3.1 O PROCESSO DE ALFABETIZAÇÃO E LETRAMENTO > O \textbf{educar} se modifica em relação ao aprender. > Significa colocar -se perante o mundo, criar \textbf{um} evento, habitar as diferentes situações.
Significa \textbf{educar} -se. > Quem participa de \textbf{um} percurso educativo de fato coloca em jogo o próprio crescimento e o faz com base nas próprias expectativas e do próprio projeto. > Há \textbf{uma} recursividade relacional entre quem \textbf{educa} e quem é \textbf{educado}, entre quem aprende e quem ensina. > Há participação, paixão, compaixão, emoção. > Carla Rinaldi \textbf{Um} dos desafios de \textbf{uma} sociedade democrática é a garantia da Educação como \textbf{um} Direito de Todos.
Essa expressão faz -nos pensar sobre a plurissignificação do vocábulo.
A palavra “ Todos ”, nesse contexto da educação, refere -se a diferentes \textbf{sujeitos} constituídos por histórias também diversas, com experiências únicas e singulares.
Esse entendimento de \textbf{sujeito} \textbf{exige} \textbf{uma} prática pedagógica \textbf{que} valorize a subjetividade como inteireza e integridade e \textbf{que} organize o processo educativo a partir do \textbf{que} já foi construído em seu percurso formativo \textbf{que} se \textbf{configura} como \textbf{constitutivo} e \textbf{constituinte} da formação humana ( \textbf{SANTA} CATARINA, 2014 ).
Ao considerarmos como parte da formação humana todas as etapas da Educação Básica, devemos nos perguntar : \textbf{Que} imagem de \textbf{humanidade} e de criança nós temos?
Quem são as crianças \textbf{que} chegam aos anos iniciais do Ensino Fundamental aos seis anos de idade ( Lei No 11.274, de 6 de fevereiro de 2006 )?
Como elas aprendem a ler e a escrever?
Como se apropriam do sistema de escrita?
Como o professor alfabetizador conduz o percurso formativo dos \textbf{sujeitos} em processo de alfabetização de modo \textbf{que} possam se apropriar da leitura e da escrita, preferencialmente no primeiro e no segundo anos, como orientado na Base Nacional Comum Curricular-BNCC ( 2017 )?
Como organizar o tempo e o espaço para acolher todas as crianças e garantir -lhes o aprendizado da leitura e da escrita em uma/na perspectiva da alfabetização e do letramento, indissociavelmente, sem desconsiderar as especificidades dos \textbf{sujeitos}?
Onde está o futuro, o \textbf{que} é o futuro?
Onde está ( ou o \textbf{que} é ) o novo?
Qual “ futuro ” é possível projetarmos juntos ( e como )?
O \textbf{que} pode parecer \textbf{uma} digressão \textbf{configura} -se, na verdade, como \textbf{uma} reflexão a respeito dos tempos, dos espaços e das práticas pedagógicas \textbf{que} são desenvolvidas com esses \textbf{sujeitos} da infância.
Desse modo, ter o entendimento da criança como \textbf{um} \textbf{sujeito} de direitos, potente, é compreender, também, a trajetória das crianças antes de chegar à escola com todas as suas especificidades ; é reconhecer a importância da Educação Infantil como \textbf{um} dos meios de potencializar a formação humana, haja vista ter \textbf{um} contexto familiar já \textbf{vivido} também, \textbf{que} tem “ [... ] o objetivo de ampliar o universo de experiências, conhecimentos e habilidades dessas crianças, diversificando e consolidando novas aprendizagens ” ( BRASIL, 2017, p. 34 ).
Nessa amplitude de experiências, os textos constituem práticas discursivas utilizadas em situações reais do cotidiano e estão presentes nas práticas sociais de leitura e de escrita.
Podemos \textbf{dizer} \textbf{que} nesse continuum do processo formativo, \textbf{que} inicia na Educação Infantil ( ou mesmo fora dela ), as crianças, gradativamente, ampliam seus repertórios linguísticos e culturais por meio de diferentes linguagens : oralidade, gestos, desenhos, brincadeiras, pinturas, esculturas, entre outras.
Assim sendo, é por meio das diferentes linguagens \textbf{que} as crianças se comunicam e expressam seu pensamento ; quanto mais forem oportunizadas vivências em \textbf{que} elas possam \textbf{viver} intensamente essas experiências, mais se desenvolvem integralmente.
Isso \textbf{implica} considerar \textbf{que}, ao ingressar no Ensino Fundamental, as crianças de seis anos \textbf{necessitam} se expressar por meio de múltiplas linguagens, e \textbf{que} as brincadeiras, a imaginação e a fantasia constituem seus modos de ser e \textbf{viver} no mundo.
Nesse sentido, compartilhamos da fala de Kramer ( 2007 ) : > Educação infantil e ensino fundamental são indissociáveis : ambos envolvem conhecimentos e afetos ; saberes e valores ; cuidados e atenção ; seriedade e riso.
O cuidado, a atenção, o acolhimento estão presentes na educação infantil ; a alegria e a brincadeira também.
E, com as práticas realizadas, as crianças aprendem.
Elas gostam de aprender.
Na educação infantil e no ensino fundamental, o objetivo é atuar com liberdade para assegurar a apropriação e a construção do conhecimento por todos.
Na educação infantil, o objetivo é garantir o acesso, [... ] a vagas em creches e pré-escolas, assegurando o direito da criança de brincar, criar, aprender.
Nos dois, temos grandes desafios : o de pensar a creche, a pré-escola e a escola como \textbf{instâncias} de formação cultural ; o de \textbf{ver} as crianças como \textbf{sujeitos} de cultura e história, \textbf{sujeitos} sociais. [... ].
Defendemos \textbf{aqui} o ponto de vista de \textbf{que} os direitos sociais precisam ser assegurados e \textbf{que} o trabalho pedagógico precisa levar em conta a \textbf{singularidade} das ações infantis e o direito à brincadeira, à produção cultural tanto na educação infantil quanto no ensino fundamental. ( KRAMER, 2007, p. 20 ).
Ao entrar na Escola²¹, a criança, ao deparar -se com os diversos símbolos, com a combinação entre eles para compor diferentes sentidos e significados, vai sentir -se motivada à apropriação da leitura e da escrita.
Nesse sentido, deve fazer parte do planejamento pedagógico a compreensão da necessidade de aprendizagem do código escrito com o objetivo maior não só da alfabetização por si só, mas da alfabetização para o exercício pleno de cidadania, iniciado ainda nos primeiros anos da escolarização do Ensino Fundamental, \textbf{uma} vez \textbf{que} “ [... ] ela [ a escrita ] condiciona a aquisição de informação na nossa sociedade e compreende a aquisição de conhecimentos e habilidades matemáticas e científicas ” ( MORAIS, 2012, p. 53 ), o \textbf{que} justifica firmarmos o \textbf{que} estamos denominando de alfabetização com e para o letramento. ²¹ É \textbf{preciso} \textbf{lembrar}, também, \textbf{que} crianças de seis e sete anos têm necessidades de movimento e \textbf{ludicidade} para se expressar e desenvolver -se cognitiva e socialmente.
Isso requer \textbf{uma} organização de tempo e de espaço \textbf{que} atenda às necessidades \textbf{inerentes} a essa fase da infância, considerando \textbf{que}, ao entrar nos anos iniciais, elas não deixam de ser crianças ( BRASIL, 2007 ).
É dentro dessa organização de tempo e de espaço adequada às especificidades das crianças \textbf{que} as \textbf{abordagens} e os \textbf{métodos} de alfabetização são materializados.
Em relação aos conceitos de alfabetização e de letramento, podemos assim defini -los : a alfabetização pode ser compreendida como \textbf{um} processo de apropriação do sistema de escrita, \textbf{que} envolve o domínio do sistema alfabético-ortográfico.
Já, em relação ao letramento, segundo a Proposta Curricular do Estado de Santa Catarina ( PCSC ) ( 2014 ), deve -se, antes de tudo, compreender o conceito de cultura escrita²², dado \textbf{que} está contido nesse conceito o de letramento, “ [... ] entendido como o conjunto de usos da escrita \textbf{que} caracteriza os diferentes grupos culturais, nas diferentes esferas da atividade humana ” ( \textbf{SANTA} CATARINA, 2014, p. 107-108 ). \textbf{Convém} destacarmos \textbf{que}, para o domínio da cultura escrita, o trabalho a ser desenvolvido deve pautar -se nas diferentes experiências interativas/trocas entre crianças e professores e crianças com o meio e nos diferentes gêneros discursivos.
Conforme pode ser \textbf{visto} na PCSC : > No aprendizado do sistema de escrita alfabética, ou mesmo antes dele, é fundamental \textbf{que} os estudantes interajam por meio da escrita em contextos sociointeracionais em \textbf{que} possam construir sentidos nas relações com o outro, mediadas pela escrita, quer o professor atue como escriba e leitor, quer as crianças já \textbf{consigam} usar a escrita de modo mais autônomo e menos heterônomo.
Importa, pois, \textbf{que} os processos de ensino considerem \textbf{que} o progressivo domínio do sistema de escrita alfabética tem de se dar nos/para os/em favor dos usos sociais da escrita, no âmbito dos gêneros do discurso. ( \textbf{SANTA} CATARINA, 2014, p. 123-124 ).
A partir dessa concepção, o trabalho com a língua em situações reais de uso, materializada no texto, sob diferentes gêneros discursivos, orais ou escritos, passa a ser a base tanto para a alfabetização quanto para o letramento.
Conforme a BNCC, o ensino da língua deve estar \textbf{ancorado} em práticas de linguagem \textbf{que} são produtos culturais \textbf{que} “ [... ] organizam e estruturam as relações humanas ” ( BRASIL, 2017, p. 60 ) ; desse modo, as atividades docentes, em se tratando também de alfabetização, devem \textbf{embasar} -se em gêneros discursivos.
Importa considerarmos também \textbf{que}, no processo de alfabetização e no decorrer do percurso formativo, a língua é mais do \textbf{que} \textbf{um} conjunto de símbolos.
Há regras de combinação entre esses símbolos, para \textbf{que} eles possam ter significados. \textbf{Convém} \textbf{lembrarmos}, ainda, \textbf{que} existem : a ) a escolha desses símbolos pelos \textbf{sujeitos}, para \textbf{que} tenham \textbf{um} determinado significado ou outro, conforme seus propósitos ( intencionalidades ) ; b ) a relação entre \textbf{sujeitos} ; c ) o contexto ; d ) a função social da língua.
Assim, é fundamental para o exercício da cidadania o trabalho com a língua viva, como prática social. ²² Modo da sociedade se organizar, tendo como fundamento a escrita, abrangendo maneiras de comportamento, “ [... ] valorações e saberes construídos pela \textbf{humanidade} e presentes na forma como os \textbf{sujeitos} se \textbf{inter-relacionam} em seu tempo e para além dele por meio dessa mesma escrita ”. ( \textbf{SANTA} CATARINA, 2014, p. 107-108 ).
O estado de Santa Catarina busca alfabetizar todas as crianças nos dois primeiros anos do Ensino Fundamental ( 1º e 2º anos ).
Isso significa \textbf{que} essas crianças devem dominar o código da escrita ( fonemas e grafemas ) e a sua função na constituição da palavra, utilizada para interagir com os mais diversos interlocutores na sociedade.
Esse processo, devido à complexidade \textbf{que} envolve o seu aprendizado, poderá dar -se a partir de diferentes \textbf{abordagens} \textbf{metodológicas}, especialmente no \textbf{que} tange à compreensão do modo como as crianças aprendem a ler e a escrever.
De acordo com Soares, > [... ] para trilhar \textbf{um} caminho, é necessário conhecer seu curso, seus meandros as dificuldades \textbf{que} se interpõem, alfabetizadores ( as ) \textbf{dependem} do conhecimento dos caminhos da criança – dos processos cognitivos e linguísticos de desenvolvimento e aprendizagem da língua escrita – para orientar seus próprios passos e os passos das crianças [... ]. ( SOARES, 2017, p. 352 ).
Nesse sentido, diante da complexidade \textbf{que} constitui a infância, pensar indicações \textbf{metodológicas} para o processo inicial de alfabetização, assim como para o percurso formativo \textbf{que} os \textbf{sujeitos} farão no decorrer do Ensino Fundamental²³, significa partir desse contexto vivenciado e apresentado pelas crianças e \textbf{que} continua as constituindo ao longo do seu percurso formativo, no Ensino Fundamental.
Isso \textbf{implica} refletir e reorganizar esse lugar, os tempos e os espaços na escola, de modo a considerar as suas especificidades, pois todo esse contexto envolve o trabalho docente. \textbf{Uma} perspectiva \textbf{que} (re)significa a relação do aprender, no sentido de \textbf{que} se considere como valor de aprendizagem algo único e fundamentado nas culturas da infância, em \textbf{que} os valores da brincadeira, da diversão, das emoções, dos sentimentos são reconhecidos como elementos essenciais para cada processo autêntico, educativo e de elaboração do conhecimento.
Esse processo de aprendizagem dá -se a partir do protagonismo dos \textbf{sujeitos} envolvidos ; assim, professores criam sua metodologia de ensino e as crianças – como interlocutoras – participam dessa criação da aula, por meio de suas falas e ações, pois são sujeitos ativos na aprendizagem.
Importa \textbf{dizer}, também, \textbf{que} essa ação ocorre em articulação com o contexto das práticas de linguagem, articulados aos diferentes campos de atuação ( Campo da vida cotidiana, Campo artístico-literário, Campo das práticas de estudo e pesquisa, Campo jornalístico-midiático e Campo de atuação na vida pública)²⁴. ²³ Vale ressaltar \textbf{que}, no processo de alfabetização e anos iniciais do Ensino Fundamental, o professor, ao planejar, \textbf{necessita} considerar, no decurso da apropriação da leitura e da escrita, o desenvolvimento da alfabetização científica, de todos os conceitos, dos componentes curriculares e dos conteúdos transversais ; é necessário, dessa maneira, \textbf{que} se considere os conteúdos de todos os componentes e áreas do conhecimento como potencializadores do planejamento docente, de forma a garantir, progressivamente, o processo de elaboração \textbf{conceitual}.
Para dar sentido e significado às práticas e às aprendizagens, o trabalho pedagógico \textbf{necessita} ser desenvolvido de modo interdisciplinar e inclusivo. ²⁴ Os campos de atuação orientam a seleção de gêneros, de práticas, de atividades e de procedimentos em cada \textbf{um} deles.
Diferentes recortes são possíveis quando se pensa em campos.
As \textbf{fronteiras} entre eles são tênues, pois reconhece -se \textbf{que} alguns gêneros incluídos em \textbf{um} determinado campo estão também referenciados a outros, existindo trânsito entre esses campos.
Práticas de leitura e de produção da escrita ou oral do campo jornalístico-midiático conectam -se às de atuação na vida pública. \textbf{Uma} reportagem científica transita tanto pelo campo jornalístico-midiático quanto pelo campo de divulgação científica ; \textbf{uma} resenha crítica pode pertencer tanto ao campo jornalístico quanto ao \textbf{literário} ou de investigação.
Enfim, os exemplos são muitos.
É \textbf{preciso} considerar, então, \textbf{que} os campos se interseccionam de diferentes maneiras. \textbf{Contudo}, o mais importante a se ter em conta e \textbf{que} justifica sua presença como organizador do componente é \textbf{que} os campos de atuação permitem considerar as práticas de linguagem – leitura e produção de textos orais e escritos – \textbf{que} neles têm lugar em \textbf{uma} perspectiva situada, o \textbf{que} significa, nesse contexto, \textbf{que} o conhecimento metalinguístico e semiótico em jogo – conhecimento sobre os gêneros, as configurações textuais e os demais níveis de análise linguística e semiótica – deve poder ser revertido para situações significativas de uso e de análise para o uso ( BRASIL, 2017, p. 85 ).
Nesse processo, buscando focalizar, objetivamente, indicações \textbf{metodológicas} para o processo de alfabetização, é \textbf{preciso} considerar quatro eixos de ação, quais sejam : oralidade, leitura, escrita e análise linguística/semiótica ( análise e reflexão sobre a língua ), articuladamente²⁵ ; além disso, pensar/planejar estratégias \textbf{que} visem à compreensão do sistema de escrita pela criança, o \textbf{que} passa pela compreensão da sua relação com a escrita e as hipóteses \textbf{que} constrói sobre a escrita, nas suas produções informais e espontâneas.
É também a partir dessa interação discursiva \textbf{que} o educador intervém com ações estratégicas \textbf{que} levem à apropriação/compreensão do sistema alfabético ( fonográfico ) da escrita.
Isso pode ser feito por meio de estratégias \textbf{que} envolvam os nomes das crianças, gêneros discursivos diversos da vida cotidiana ( rótulos, listas, cantigas folclóricas, diferentes narrativas, gêneros orais e escritos, entre outros ) e \textbf{que} agrupem determinados campos semânticos, de modo \textbf{que} seja possível criar condições para o desenvolvimento da consciência fonológica, análise e reflexão sobre a escrita a partir dos seus usos e práticas discursivas²⁶, mas sempre com ações articuladas entre os diferentes eixos \textbf{que} envolvem o planejamento no processo de alfabetização ( oralidade, leitura, escrita e análise linguística/semiótica ) e componentes curriculares diversos.
Entendendo \textbf{que} é nesse contexto \textbf{que} a criança se apropria do sistema de escrita, temos a responsabilidade formal pela \textbf{sistematização} da alfabetização/do sistema alfabético ( fonográfico ) da escrita, nos primeiros anos do Ensino Fundamental.
Para tanto, é necessário planejar ações sistemáticas e \textbf{que} promovam esse aprendizado por meio das diferentes práticas de linguagem.
Tais ações envolvem : ²⁵ Na leitura do Quadro apresentado ao final deste texto ( Apêndice A ), é possível perceber \textbf{que} os conceitos/conteúdos e as habilidades, por organização didática, estão separados, mas eles precisam ser lidos de modo articulado ; desse modo recomenda -se \textbf{que} se comece pelos objetos de conhecimentos, a habilidade \textbf{que} deve ser desenvolvida e depois os conteúdos como meio para desenvolver as habilidades.
No caso do componente curricular língua portuguesa, por exemplo, a leitura dos eixos oralidade, leitura, escrita e análise linguística também devem ser lidos de modo articulado, pois não trabalhamos o desenvolvimento desses conceitos em separado, mas no conjunto do planejamento. ²⁶ Leal, Albuquerque e Rios ( 2005 apud BRASIL, 2007 ) citam algumas brincadeiras \textbf{que} fazem parte da nossa cultura e envolvem a linguagem, \textbf{que} podem ajudar no processo de alfabetização e auxiliar os alfabetizandos na automatização dos valores dos grafemas : cantar músicas e cantigas de roda, recitar parlendas, poemas, quadrinhas, adivinhas, jogo da forca, entre outras.
As \textbf{autoras} ressaltam \textbf{que} os jogos fonológicos – aqueles \textbf{que} dirigem a atenção da criança para as semelhanças e diferenças sonoras entre as palavras – podem ser \textbf{poderosos} aliados dos professores no ensino da leitura.
No caso da apropriação do sistema alfabético, tais jogos possibilitam à criança “ [... ] manipular as unidades sonoras/gráficas ( palavras, sílabas, palavras ), a comparar palavras ou partes delas [... ] ” ( BRASIL, 2007, p. 81 ), a usar pistas para ler e escrever outras palavras, avançando, dessa forma, na aprendizagem inicial da leitura. i.
Situações de uso real das práticas discursivas ( organizar ambiente alfabetizador \textbf{que} garanta a circulação de diferentes suportes de gêneros e práticas discursivas - literatura infantil, jornais, músicas, rótulos, cartazes, placas, jornais, bilhetes, avisos, crachás, recados, revistas, entre outros gêneros escritos e da oralidade ). ii.
Jogos e atividades lúdicas diversas em \textbf{que} as crianças sejam envolvidas/desafiadas a comparar e relacionar palavras entre si, com suas ilustrações, etc. iii.
Atividades em \textbf{que} o \textbf{aprendiz} da escrita é desafiado a produzir escrita espontânea, completar textos \textbf{conhecidos} de diferentes modos, relacionar fonema/grafema, circular palavras \textbf{conhecidas}, fazer associações/comparações tanto na escrita quanto nos efeitos de sentido \textbf{que} esta produz em seus interlocutores, etc. iv.
Alfabeto móvel : é \textbf{desejável} \textbf{que} todas as crianças possam ter o seu alfabeto móvel ( de preferência colorido ) para \textbf{que}, em grupos e individualmente, possam “ exercitar ” diferentes possibilidades de produção de palavras e textos. v.
Diferentes experiências com gêneros \textbf{literários} diversos, cotidianamente, de modo \textbf{que} eles possam ampliar o universo vocabular, compreensão e leitura de mundo, reflexões sobre diferentes temas e conceitos. \textbf{vi}.
Planejamento docente, visando o desenvolvimento do pensamento complexo, envolvendo os diferentes eixos e componentes curriculares, ( interdisciplinar ) articulado aos diferentes campos de atuação.
A investigação, a pesquisa, ( claro, adequada a essa etapa de ensino ), deve \textbf{permear} todo o percurso formativo, quando do planejamento de diferentes situações desencadeadoras de aprendizagem²⁷ ( \textbf{SANTA} CATARINA, 2014 ).
É importante \textbf{que} haja \textbf{um} trabalho contínuo e sistemático por parte dos docentes com gêneros discursivos diversos, o \textbf{que} deve ocorrer de modo articulado/simultâneo durante o processo de apropriação do código ortográfico e durante o percurso formativo, em \textbf{um} continuum progressivo das aprendizagens no qual o professor amplia o nível de complexidade dos conceitos, pela análise linguística e de reflexão sobre a escrita, abordando a ortografia “ correta ” das palavras²⁸.
Isso, no geral, ocorre com mais ênfase a partir do segundo ano, quando os \textbf{aprendizes} já dominam o código alfabético da nossa escrita, mas é \textbf{um} trabalho \textbf{que} se segue pelos próximos anos do Ensino Fundamental, nos quais se deve garantir a diminuição de \textbf{erros} ortográficos ao mesmo tempo \textbf{que} se amplia a expressão escrita.
Para aprender a escrever, as crianças têm \textbf{um} longo processo até entender \textbf{que} a escrita representa a fala, porém essa representação não é direta.
Quando aprendem a escrever em geral, ainda escrevem como falam e, consequentemente, grafam do modo como relacionam os sons e as letras \textbf{que} \textbf{necessitam} para escrever tais fonemas.
Respeitar esse processo é \textbf{um} avanço no aprendizado da escrita, pois favorece a manifestação oral e espontânea da criança, dando liberdade para expor suas ideias, tanto na oralidade quanto por meio da escrita, sem ser cerceada pela “ correção ” \textbf{uma} vez \textbf{que} se entende essa fase como processo de compreensão do sistema de escrita, de valorização do discurso oral e escrito ( e não menosprezo em detrimento da forma “ correta ” de escrita ).
É \textbf{preciso} não só respeitar, mas também intervir²⁹. ²⁷ Nesse sentido, é importante \textbf{lembrar} \textbf{que} : “ Dialogar com as diferentes formas do conhecimento \textbf{exige} pensar em estratégias \textbf{metodológicas} \textbf{que} permitam aos estudantes da Educação Básica desenvolver formas de pensamento \textbf{que} lhes possibilitem a apropriação, a compreensão e a produção de novos conhecimentos.
Tais estratégias nos \textbf{remetem} à compreensão da atividade orientadora de ensino ” ( \textbf{SANTA} CATARINA, 2014, p. 157 ). ²⁸ Importante \textbf{lembrar}, também, \textbf{que}, no processo de alfabetização, tratamos o “ \textbf{erro} ortográfico ” sempre como processo, por isso a necessidade de acompanhamento individualizado constante, de modo \textbf{que} as intervenções possam ocorrer sistematicamente e levando a reflexão/comparação entre os diferentes fonemas e modos de grafá -los, para \textbf{que} progressivamente eles possam incorporar os diferentes fonemas e seus diferentes modos de escrever. ²⁹ Vale ressaltar \textbf{que} “ [... ] não podemos [... ] esperar \textbf{que} eles [ os \textbf{aprendizes} da escrita ] aprendam ortografia apenas com o tempo ou sozinhos.
É \textbf{preciso} garantir \textbf{que}, \textbf{enquanto} avançam em sua capacidade de produzir textos, \textbf{vivam} simultaneamente oportunidades de registrá -los cada vez mais de forma correta ” ( MORAIS, 2000, p. 22-23 ).
Em relação ao processo de avaliação na alfabetização, ele também está \textbf{ancorado} nas concepções \textbf{aqui} elencadas, as quais devem subsidiar/retroalimentar a prática docente continuamente.
Estamos falando de \textbf{uma} concepção de avaliação diagnóstica, formativa, processual, contínua e sistemática.
Nessa concepção, o docente pode acompanhar o processo ensino-aprendizagem do \textbf{sujeito} e, a partir desse acompanhamento, retomar o \textbf{que} não foi aprendido, (re)planejar, de modo a garantir as aprendizagens, na perspectiva da formação integral³⁰.
Desse modo, a “ [... ] avaliação deve servir como \textbf{um} instrumento de inclusão e não de classificação e/ou \textbf{exclusão}.
Deve ser \textbf{um} indicador não apenas do nível de desenvolvimento do estudante, como também das estratégias pedagógicas e das escolhas \textbf{metodológicas} do professor ” ( \textbf{SANTA} CATARINA, 2014, p. 46 ). ³⁰ Dado seu caráter formativo, contempla pelo menos três etapas : a de diagnóstico, a de intervenção e a de replanejamento.
O trabalho de diagnóstico ocorre quando o professor verifica a aprendizagem \textbf{que} o estudante realizou ou não, compreendendo as possibilidades e as dificuldades do processo, no momento.
A intervenção se dá quando o professor retoma o percurso formativo, após constatar \textbf{que} não houve suficiente elaboração \textbf{conceitual}, e, por isso, reorganiza o processo de ensino possibilitando ao sujeito novas oportunidades de aprendizagem.
O replanejamento é \textbf{uma} \textbf{tarefa} \textbf{que} se faz necessária sempre \textbf{que} as atividades, estratégias de ensino e seus respectivos resultados não se evidenciarem suficientes.
Ao longo do desenvolvimento das três etapas, é fundamental \textbf{que} se considere a \textbf{sistematização}, a elaboração e a apropriação de conhecimentos, na forma de registros, relatos e outros instrumentos como subsídios para a avaliação.
Neste âmbito, importa \textbf{que} os registros considerem relatos dos \textbf{sujeitos} acerca das suas próprias atividades, sejam elas práticas, \textbf{teóricas} ou lúdicas, bem como outros instrumentos \textbf{que} subsidiem a avaliação ( \textbf{SANTA} CATARINA, 2014, p. 47 ).
Por fim, é importante reiterar \textbf{que} o processo de alfabetização e letramento, em \textbf{uma} perspectiva mais ampla, ocorre ao longo do percurso formativo e precisa ser \textbf{compromisso} de todas as áreas e de todos os componentes curriculares ; dessa maneira, todos devem trabalhar considerando o texto como articulador da prática pedagógica, os diferentes gêneros discursivos como estratégia de ensino, como meio para elaborar suas sínteses.
Nessa lógica, entende -se \textbf{que} quanto mais os \textbf{sujeitos} ampliam suas aprendizagens, elaboram e reelaboram conhecimentos/desenvolvem o pensamento complexo, mais alfabetizados e letrados eles se tornam.
A partir dessas considerações, esperamos ter auxiliado os(as) educadores(as) catarinenses com reflexões sobre o percurso a ser desenvolvido no processo de alfabetização.
No entanto, são apenas algumas possibilidades, o caminho a ser trilhado \textbf{necessita} ser elaborado de forma a considerar a historicidade e os contextos dos \textbf{sujeitos} envolvidos, suas especificidades humanas.
Cabe aos educadores(as), em \textbf{um} processo contínuo de estudo e de formação continuada, buscarem o domínio desses conceitos³¹ e, como \textbf{autores} da sua própria prática pedagógica, galgarem os seus percursos a partir da realidade em \textbf{que} se encontram imersos.
Desse modo, desejamos \textbf{que} as escolas possam cada vez mais ser o lugar da cultura mais elaborada ( MELLO ; FARIAS, 2010 ) e \textbf{que} todos(as) os(as) educadores(as), possam enfrentar esse desafio, de maneira a vivenciarem, cotidianamente, experiências pedagógicas \textbf{que} possam garantir a todos(as) os(as) estudantes catarinenses o direito de aprender a ler e a escrever, de alfabetizar -se, no âmbito da alfabetização e do letramento, de modo indissociável e progressivamente, pleno.

% matched lemmas: abordagem, ancorar, aprendiz, aqui, autor, compromisso, conceitual, configurar, conhecido, conseguir, constituinte, constitutivo, contudo, convir, depender, desejável, dizer, educar, embasar, enquanto, erro, exclusão, exigir, fronteira, humanidade, implicar, inerente, instância, inter-relacionar, lembrar, literário, ludicidade, metodológico, método, necessitar, permear, poderoso, preciso, que, remeter, santo, singularidade, sistematização, sujeito, tarefa, teórico, um, ver, viver
\end{textsample}
